\noindent\hspace*{1.2em}\kbvlabel{ಉದಾ}\hspace{0.6em}ಪುಷ್ಪ ಸಮುದಾಯದಲ್ಲಿ ಗುಲಾಬಿಯು ಹೇಗೋ ಕವಿ ಸಮುದಾಯದಲ್ಲಿ ಲಕ್ಷ್ಮೀಶನು ಹಾಗೆ.\par

\kbvsechead{(೨)\hspace{0.45em}ರೂಪಕಾಲಂಕಾರ :}

\indent “ಸೀತಮ್ಮನವರ ಸೊಸೆ ಮಹಾಲಕ್ಷ್ಮಿಯೇ ಸರಿ.”—ಈ ವಾಕ್ಯದಲ್ಲಿ, ಸೀತಮ್ಮನವರ ಸೊಸೆಯೇ ಮಹಾಲಕ್ಷ್ಮಿ ಎಂಬ ಅರ್ಥ ಬರುತ್ತದೆ. ಅಂದರೆ ಸೀತಮ್ಮನವರ ಸೊಸೆಗೂ ಮಹಾಲಕ್ಷ್ಮಿಗೂ ಯಾವುದೇ ಭೇದವಿರುವುದಾಗಿ ತಿಳಿದುಬರುವುದಿಲ್ಲ... ಉಪಮಾನ ಉಪಮೇಯಗಳೆರಡೂ ಒಂದೇ ಎಂಬ ಅಭೇದರೂಪಕ ಇಲ್ಲಿ ಕಂಡುಬರುತ್ತದೆ. ಈ ಬಗೆಯ ಅರ್ಥಾಲಂಕಾರವನ್ನು ರೂಪಕಾಲಂಕಾರವೆನ್ನುವರು.\par

\noindent\hspace*{1.2em}\kbvlabel{ಉದಾ}\hspace{0.6em}ಗಾನವಿಶಾರದೆ ಶ್ರೀಮತಿ ಸುಬ್ಬಲಕ್ಷ್ಮಿಯು ಸಾಕ್ಷಾತ್‌ ಸರಸ್ವತಿ.\par

\kbvsechead{(೩)\hspace{0.45em}ಉತ್ಪ್ರೇಕ್ಷಾಲಂಕಾರ :}

\indent “ಹೈನುಗಾರಿಕೆಗೆ ಪ್ರೋತ್ಸಾಹ ಸಿಕ್ಕಿದ ಕಾರಣ ಕರ್ನಾಟಕದ ರಾಜ್ಯದಲ್ಲಿ ಹಾಲಿನ ಹೊಳೆ ಹರಿಯುತ್ತಿದೆ.”—ಈ ವಾಕ್ಯದಲ್ಲಿ ಉತ್ಪ್ರೇಕ್ಷೆ ಎದ್ದು ಕಾಣುತ್ತಿದೆ. ರಾಜ್ಯದಲ್ಲಿ ಹಾಲಿನ ಉತ್ಪಾದನೆ ಗಣನೀಯವಾಗಿ ಹೆಚ್ಚಿದೆಯೆಂಬುದನ್ನು ತಿಳಿಸುವ ಉದ್ದೇಶವಿಟ್ಟುಕೊಂಡು ವಸ್ತುಸ್ಥಿತಿಯನ್ನು ಅತಿಶಯವಾಗಿ ಬಣ್ಣಿಸುವುದು ಇಲ್ಲಿನ ವೈಶಿಷ್ಟ್ಯ... ಈ ಬಗೆಯ ಅರ್ಥಾಲಂಕಾರವನ್ನು ಉತ್ಪ್ರೇಕ್ಷಾಲಂಕಾರವೆನ್ನುವರು.\par

\noindent\hspace*{1.2em}\kbvlabel{ಉದಾ}\hspace{0.6em}ಭೀಮನು ಶಿಶುವಾಗಿದ್ದಾಗ ತಾಯ ತೊಡೆಯಿಂದ ಕೆಳಗುರುಳಿ ಬಿದ್ದುದರಿಂದ ಆ ಅಡಿಗಲ್ಲೇ ಚೂರ್ಣವಾಯಿತಂತೆ !\par

\kbvsechead{(೪)\hspace{0.45em}ದೃಷ್ಟಾಂತಾಲಂಕಾರ :}

\indent “ಮನೆಗೊಂದು ಮುದುಕಿ, ಒಲೆಗೊಂದು ಕೊರಡು.”—ಈ ಗಾದೆ ಮಾತಿನಲ್ಲಿ ಪ್ರಧಾನವಾದ ಅಂಶವು ಪೂರ್ವಾರ್ಧದಲ್ಲಿ ಅಡಕವಾಗಿದೆ. “ಮನೆ ಗೊಂದು ಮುದುಕಿ” ಅವಶ್ಯಕವಾಗಿ ಇರಬೇಕು. ಆ ಮುದುಕಿಯು ತನ್ನ ಅನುಭವಗಳನ್ನು ಮುಂದಿನ ಪೀಳಿಗೆಗೆ ಧಾರೆ ಎರೆಯುವ ಮೂಲಕ ಆ ಮನೆಯ\kbvcontinues
